% \section{Introduction and Module Walkthrough}

\begin{frame}{Portfolio Exam}
    \mytext{
    \begin{itemize}
        \item 50\% presentation of an economy, 20 minutes
        \item 50\% written exam at the end of term
    \end{itemize}
    }
    \mycomment{
        The exam is portfolio. That means it is a unitary exam, and the final grades are composed in view of all results of all parts and of  students at the end. 

        That means also that there will not be a publication of exact grades for any part of the portfolio.

        This also means that it is not possible to repeat a single part of the portfolio -- you can only repeat the whole exam with all parts (and the entire module) if you fail, or upon proof of illness or other force majeure.

        For all other issues regarding the exam: note that you have to deal with your Prüfungsamt and Prüfungsausschuss. The lecturer is responsible only for matters of content. The Prüfungsamt also sets the date of the written exam.
    }
\end{frame}

\begin{frame}{Presentation}
    \mytext{
    Choose an Economy: a \textbf{country, state or region} (\enquote{Gebietskörperschaft})!

    Examples:
    
    Bremen (Germany); Oaxaca (Mexico); Maharashtra (in India); Hainan (in China); Arkansas (the US); Quebec (Canada); KwaZulu-Natal (South Africa); The Netherlands; Botswana; etc.

    Someone should please cover Japan, or a part of Japan.

    }
    \mycomment{
        Rather than holding your presentation live in person, I want you to present on video. We will discuss technical options. The result shall be a video of 20 minutes. 
    }
\end{frame}

\begin{frame}{Presentation}
    \mytext{    
    Your presentation has two parts; part one -- you present  economic data; part two -- you present the  governments economic policy.
    
    \begin{itemize}
        \item 20 minutes, no fancy animations
        \item using HSB style guide / slide layout
        \item (live) recorded video
        \item students will give feedback and grade each other
        \item[] that feeds into the grading.
    \end{itemize}
    }
    \mycomment{
    
        To receive a grade, you yourself must also comment and grade a minium number of other videos. These mutual gradings (kept anonymous) will form part of the final grade.
        This will allow all students to scrutinise your performance and hence make grading more transparent.
    }
\end{frame}

\begin{frame}{Presentation}
    \mytext{
    Part I: Present core economic data: present, history and trends
        \begin{itemize}
            \item core geo- and demographic data
            \item constitution, institutions
            \item GDP, per head, inequality
            \item inflation, currency
            \item employment, education
            \item core industries / prices
            \item import / export dependency
            \item trade openness
        \end{itemize}
    }
    \mycomment{
    Use a \enquote{zoom-in} or \enquote{focus} technology. Present the country data as framework, and present the largest components.
    
    When looking at a large state that is organised as a federal state (USA, China, Germany, India, Brazil etc.), zoom into a sub-unit (federal unit).

    When discussing the industry or core industries, focus on the most important one. For instance, Germany is a highly industrialised country; the most important ones are car manufacturing, machines, and chemistry; for chemistry, these are the latest numbers... etc.

    You can already come to some conclusion at the end of this part, as to where the main challenges lie.
    }

\end{frame}

\begin{frame}{Presentation}
    \mytext{
    Part II: Economic Policy
        \begin{itemize}
            \item stated objectives, real objectives
            \item investment subsidies
            \item production subsidies
            \item regulation and standards
            \item wages, social policy, education
            \item export subsidies
            \item tax rebates, special economic areas
            \item government procurement
            \item innovation policy
            \item monetary policy 
        \end{itemize}
    }
    \mycomment{
        In the second part, you explain the current governments economic policy. You must demonstrate critical understanding. How do you view these policies?
        
        Finally, you should also mention: What are the main opposition party proposals? What do they criticise an what do they propose instead? What does the international economic community, the IMF and the World Bank, propose?    
    }
\end{frame}
